\chapter{Практика}

\section{Построение кластеризации}

Построим модель кластеризации и визуализируем результат


\begin{figure}[H]
    \centering
    \includegraphics{pca.png}
    \caption{Визуализаци кластеризации}
\end{figure}

Можем увидеть, что модель разбила семена на 3 группы (кластера). \\

\section*{Интерпретация silhouette plot}

Также для анализа качества кластеризации используется \textbf{silhouette plot}. \\

\begin{figure}[H]
    \centering
    \includegraphics{splt.png}
    \caption{Качество кластеризации}
\end{figure}



Данный график представляет собой \textbf{silhouette plot}, используемый для оценки качества кластеризации.

\begin{enumerate}
    \item \textbf{Ось $x$} — значение silhouette coefficient для каждого объекта:
    \begin{itemize}
        \item $s \approx 1$ — объект хорошо принадлежит своему кластеру;
        \item $s \approx 0$ — объект находится на границе между кластерами;
        \item $s < 0$ — возможная неправильная кластеризация.
    \end{itemize}

    \item \textbf{Ось $y$} — номера кластеров. Каждый цветной блок соответствует отдельному кластеру.

    \item \textbf{Синяя вертикальная линия} — среднее значение silhouette coefficient по всей выборке.

    \item \textbf{Интерпретация результата}:
    \begin{itemize}
        \item среднее значение около $0.4$ указывает на \textit{умеренно хорошее} качество кластеризации;
        \item большинство значений положительные, что говорит о корректном разделении кластеров;
        \item наличие значений ближе к $0$ указывает на пересечение границ кластеров.
    \end{itemize}
\end{enumerate}


\section*{Метод «локтя» (Elbow Method)}

Метод «локтя» используется для выбора оптимального числа кластеров $k$ в алгоритмах кластеризации, таких как $K$-means.


\begin{figure}[H]
    \centering
    \includegraphics{el.png}
    \caption{Метод локтя}
\end{figure}


На графике по оси $x$ отложено число кластеров $k$, а по оси $y$ — значение инерции, которое определяется как сумма квадратов расстояний между объектами и центрами их кластеров:

\[
J(k) = \sum_{i=1}^{k} \sum_{x \in C_i} \|x - \mu_i\|^2
\]

где $C_i$ — множество объектов в $i$-м кластере, а $\mu_i$ — центр соответствующего кластера.

С увеличением числа кластеров значение инерции монотонно уменьшается, однако после некоторого значения $k$ улучшение качества разбиения становится незначительным. Эта точка называется «локтем».

По представленному графику видно, что резкое снижение инерции происходит при $k = 1$ до $k = 3$, после чего уменьшение становится более плавным. Следовательно, оптимальное число кластеров находится в области $k \approx 3$ (возможно $k = 4$).

\endinput